\chapter{离线生成动力学：梦境、记忆回放与模型校准}
\label{chap:physics_of_dreams}
梦境涉及外部输入降低、内部生成增强、记忆重放、情绪加工和元认知改变等多种过程，不能由单一“升温”或“流形自由演化”概括。我们在本章把梦境作为一种受生理状态约束的离线生成特例：系统仍接收内感受和残余外部信号，行动通道受到门控，内部模型以不同于清醒任务的方式采样与组合。HSF--HD负责描述状态、输入、结构、控制和报告如何改变；神经机制与临床现象须服从相应领域证据，AgentOS的离线维护只能作为工程对照，不能反向证明人脑按同一模块运行。每项功能解释还要与无梦睡眠、清醒想象和离线训练等替代条件比较。

\section{输入与行动边界的重配置：从现实闭环到内源生成}
\label{sec:section_3039}
清醒与睡眠的差别不是从狄利克雷边界突然切换到诺伊曼边界，更不是从开放系统变成绝热封闭系统。睡眠中视觉、听觉和运动输出的增益会变化，但内感受、姿态、呼吸、温度、噪声和触觉仍可进入；梦境也会受近期事件与身体状态影响。我们把边界写成多通道增益与权限的联合状态，而不是把“现实”设为场在边界上的固定值。

令外部观测通道为 $u_k^{ext}\in\mathbb R^{d_e}$，内感受通道为 $u_k^{int}\in\mathbb R^{d_i}$，内部生成候选为 $z_k\in\mathbb R^{d_z}$，动作命令为 $a_k$。睡眠模式 $m_k$ 决定感知增益 $G_e(m_k)$、内感受增益 $G_i(m_k)$、动作门控 $G_a(m_k)$ 与生成增益 $G_z(m_k)$。工作态离散更新可写为
\[
h_{k+1}=F\!\left(h_k,G_e(m_k)u_k^{ext},G_i(m_k)u_k^{int},G_z(m_k)z_k,\Theta_k;\Delta t_k\right),
\qquad
\tilde a_k=G_a(m_k)a_k .
\]
各增益是有单位接口上的线性或非线性算子，不能把它们相加为一个抽象温度。模型还需声明采样步长、通道延迟、噪声和模式切换条件；若模式在一次步长内变化，应使用事件重置而非假定连续可导。

入睡后外部校正通常减弱，内部模型相对主导，于是生成内容更容易被当前系统当作场景证据。但“相对主导”不等于因果律失效。梦中对象仍由既有记忆、生成模型、情绪和身体信号约束；门铃、疼痛或呼吸困难也可能被编入梦境并触发醒来。系统应区分真实外部消息、内部重放、模型补全和梦中再解释，若醒后报告无法恢复来源，就只能将内容标记为低来源置信的体验报告，而不能据此重写外部事实。

感知边界的重配置可以用校准比描述。给定任务读出 $R_q$，定义外部证据敏感度
\[
\chi_e=\left\|\frac{\partial R_q(h_{k+1})}{\partial u_k^{ext}}\right\|,
\qquad
\chi_z=\left\|\frac{\partial R_q(h_{k+1})}{\partial z_k}\right\|.
\]
当 $\chi_z/\left(\chi_e+\epsilon\right)$ 增大时，报告更受内部生成影响；该比值依赖查询、通道和状态，不是梦境强度的普适标量。若外部刺激足够强，$\chi_e$ 仍可能上升并改变模式。验证需要在不同睡眠阶段施加受控声音、气味或触觉，比较报告、神经指标和觉醒概率，而不是从方程名称直接推出现象。

动作边界同样是门控而非绝对切断。快速眼动睡眠中的肌张力抑制能降低多数随意动作输出，但眼动、呼吸和局部肌肉活动仍存在；不同睡眠状态和个体也有差异。工程系统可通过影子执行器实现类似隔离：计划器产生动作候选，模拟环境返回虚拟回执，真实执行器保持只读或需要额外授权。此时必须给虚拟回执标记来源，防止离线模拟被写成真实经历。若动作门控异常开启，系统还要有硬件边界和紧急停止，不能只依赖上层意图。

边界变化会改变可辨识性。外部反馈减少时，多种内部模型可能生成相似梦境，单次报告无法判断是哪一机制主导；醒后叙述又受到遗忘和语言重构影响。我们因此把梦中状态、觉醒瞬间记录、延迟报告和后来解释分层保存。对人类研究，主观报告、行为、生理信号和神经数据各提供有限证据；对AgentOS离线维护，则可直接记录生成轨迹、模型版本、外部输入与执行门控，但即便日志完整，也不能由功能相似推出主体体验相同。

本节保留“从处理现实转向模拟现实”的直觉，但将其限定为输入—生成权重与行动权限的重新配置。验收使用残余外部刺激、内感受扰动、模拟回执污染、动作门控故障和模式快速切换。系统应识别消息来源，避免把模拟写成事实，在外部风险上升时能够唤醒或升级控制，并在恢复清醒闭环后重新校准环境版本。报告外部敏感度、生成敏感度、来源误判、动作泄漏、唤醒延迟和实际任务结果；这些指标支持离线生成模型，不证明睡眠是绝热演化或思维波在边界全反射。

\section{元控制降增益与叙事一致性的动态松弛}
\label{sec:section_4656}
原章把快速眼动睡眠概括为前额叶“沉默”、意志势能消失和逻辑度量液化，这个描述把脑区平均活动、元认知能力和叙事约束混成了一个开关。我们改用可检验的控制变量描述：某些睡眠阶段中，现实核验、冲突监测、计划保持和来源判断的有效增益可能下降，但没有理由假定全部高层调节归零；梦者仍可回避威胁、延续目标、产生局部推理，清明梦更表明元认知可以部分恢复。神经活动的区域差异是经验问题，HSF--HD只刻画功能状态之间如何耦合，不把某脑区等同于唯一的宏观控制器。

令生成状态为 $h_k\in\mathbb R^{d_h}$，情景上下文为 $e_k\in\mathbb R^{d_e}$，结构参数为 $\Theta_k\in\mathbb R^{d_\theta}$。我们另外定义元控制状态 $c_k\in[0,1]^{d_c}$、来源标签分布 $s_k\in\Delta^{n_s-1}$ 与叙事约束集合 $\mathcal C_k$。这里 $c_k$ 的分量分别控制现实核验、时间定位、身份绑定、反事实标注和动作许可；它不是意识强度，也不是焦耳能量。内部候选 $z_k$ 经过约束读出器形成下一状态：
\[
h_{k+1}=F(h_k,e_k,z_k,\Theta_k)-B(c_k)\nabla_h V(h_k;\mathcal C_k)+\eta_k,
\qquad s_{k+1}=S(s_k,h_{k+1},c_k).
\]
$V$ 是违反当前叙事约束的计算代价，$B(c_k)$ 是有界控制增益，$\eta_k$ 是模型未解释扰动。只有在 $V$ 可微、$B$ 半正定且目标与约束在该步内不变时，第二项才倾向于降低局部违约代价；它不保证梦境符合外部事实，更不保证任务成功。若人物身份、场景和时间规则在生成过程中切换，$\mathcal C_k$ 也必须同步更新，不能把所有突变解释为“逻辑距离融化”。

这一模型能区分三类常见现象。第一类是局部连贯而全局跳跃：每个片段内部满足一组约束，但片段切换时身份绑定或地点索引被重置，于是房间可以通向大海、一个人可以被另一身份替换。第二类是来源置信不足：内部生成内容被赋予过高的外部来源概率，梦者便暂时不质疑荒诞事件。第三类是元认知再入：当现实核验和自我状态读出的部分增益上升，系统可能识别“我正在做梦”，却未必获得完整场景控制。清明梦因此是多变量混合状态，不是一个宏观层突然全部上线的“上帝模式”。即使梦者能改变飞行或场景，也只是内部生成器对意图条件的响应，不能说明意志直接修改了物理世界的度量张量。

在AgentOS工程实例中，$c_k$ 可由TCE维护，SPC把当前生成内容、来源置信和自我状态压缩后重新送入 $h(t)$，CCM监测候选分支数量和一致性检查的开销。离线模式可以降低外部任务优先级、放宽联想阈值，同时保持Boundary的执行禁令和Me写入权限。这里SPC、TCE、CCM只是一个实现族：其他智能系统可以用验证器、事务控制或分层规划器实现同等功能。关键不是复制人脑睡眠，而是把“生成自由度增加”与“事实写入、结构更新和真实执行仍受约束”分开。

离散实现还需防止增益切换造成振荡。设 $c_{k+1}=\Pi_{[0,1]^{d_c}}(c_k+\Delta t_k f_c-\rho(c_k-c_{k-1}))$，其中投影保持范围，$\rho$ 抑制相邻步的剧烈反转。若检测到外部告警、呼吸异常、动作泄漏或高风险内容，控制器不应继续追求叙事自由，而应提高唤醒优先级或终止离线周期。验收场景包括普通梦、清明梦、清醒想象与无梦报告；比较片段一致性、来源判断、元认知命中、意图可控度和觉醒延迟。由于醒后报告会重构内容，未报告不等于未发生，报告荒诞也不等于当时完全没有控制。本节由此保留“意志减弱与逻辑松弛”的现象问题，却把它重建为可分解、可测量且允许反例的元控制动力学。

\section{离线采样、模型校准与有限泛化}
\label{sec:section_5559}
原章用升温、越过势垒、逆里奇流和平滑曲率解释梦境功能，容易把探索随机性、神经调质、训练损失和热力学温度视为同一量。我们只把“退火”保留为算法类比：离线状态可能扩大内部采样范围、重组远距离记忆并暴露模型冲突，但梦境是否承担这种功能必须由对照实验确认。对人类而言，睡眠对记忆和学习的重要性不等于梦的主观内容本身就是优化程序；对工程智能而言，离线训练是否有益取决于数据、目标、更新规则与回归验证，不能断言所有长期运行的系统都必须做梦。

为明确对象，设结构模型参数为 $\Theta\in\mathbb R^{d_\theta}$，近期经验缓冲为 $E_r$，长期保留集为 $E_a$，内部生成器为 $q_\phi(x\mid E_r,\Theta)$，实际任务损失为 $\ell(\Theta;x)$。离线批次由真实回放、生成变体和反事实样本组成：
\[
\mathcal B_k=\mathcal B_k^{replay}\cup\mathcal B_k^{gen}\cup\mathcal B_k^{cf},
\quad
J_k(\Theta)=\widehat L_{\mathcal B_k}(\Theta)+\lambda_a\widehat L_{E_a}(\Theta)
+\lambda_s D(\Theta,\Theta_k)+\lambda_u U(\Theta;\mathcal B_k).
\]
第一项拟合离线批次，第二项保护旧能力，第三项限制参数漂移，第四项惩罚不确定条件下的过度承诺；这些项都是计算代价，单位由实现定义，并非热量。若采用一步梯度更新 $\Theta_{k+1}=\Theta_k-\alpha_k\widehat\nabla J_k$，只有在梯度估计、光滑性和步长满足相应条件时，才可讨论训练目标的期望下降。即便 $J_k$ 下降，真实环境中的任务表现仍可能恶化，因为生成样本偏置、保留集缺失或目标错配会产生模型自我污染。

所谓“提高温度”在这里被拆为采样分布参数。若候选得分为 $r_i$，可以定义 $p_\tau(i)\propto\exp(r_i/\tau)$；增大无量纲参数 $\tau$ 会使给定候选集上的选择更平坦，却不自动产生有价值的远距离关联。候选集本身受Me、$E$、$\Theta$和检索策略限制，低质量记忆在高探索下也会被放大。反过来，低温采样并不等于清醒，确定性的搜索也能产生虚构。我们因此分别记录采样熵、语义跨度、模型不确定度、更新幅度和硬件能耗，禁止把它们合称为“认知能量”。

泛化必须通过外部留出任务验证。设更新前后在互不参与训练的任务集合 $\mathcal T_{hold}$ 上得分为 $Q(\Theta_k)$ 与 $Q(\Theta_{k+1})$，离线更新只有在新能力、旧能力、安全约束和校准误差的联合门槛都通过时才被接受。单一平均分会掩盖灾难性遗忘，因此还需报告最差任务下降、群体差异和长尾失败。若候选更新未通过，就回滚参数而保留诊断日志；如果多轮回滚集中于同一经验簇，系统应修改生成器或数据选择，而不是继续“升温”。这比假定逆里奇流天然平滑知识结构更可执行，因为参数空间没有预先给定的黎曼曲率，也没有证据说明正号曲率流能实现所需的遗忘—保留权衡。

对人类梦境，我们能提出但不能在本章宣告成立的经验假说：离线生成的片段跨度可能与随后联想任务的迁移相关，近期未解问题可能改变梦境内容，情绪显著材料可能提高被重组的概率。反例同样重要：大量梦境没有可见的问题求解价值，睡眠获益可以在没有可报告梦境时出现，过强噩梦反而损害休息和次日功能。实验应比较睡眠阶段、唤醒时点、报告内容与次日表现，并控制期待效应和回忆偏差。本节因此把“认知退火”限定为候选生成与验证机制：探索提供可能性，保留集和外部测试决定哪些变化能够提交，任何全局优化结论都需要额外条件。

\section{记忆回放、结构提交与情绪再估计}
\label{sec:section_9179}
梦境与记忆的关系不能简化为质流冲刷皮层、把经历永久刻进度量张量。睡眠中的回放、重激活和巩固是可研究的神经与行为现象，但主观梦境、海马—皮层交互和长期记忆改善不是一一对应关系；不同睡眠阶段、记忆类型和任务具有不同机制。我们在一般计算层只描述经验记录如何被选择、重建、评估并候选性地提交到结构记忆，不把某种脉冲模式直接翻译成HSF--HD方程，也不把AgentOS缓存等同于海马体。

设情景记忆 $E$ 由带版本的事件记录 $e_j=(o_j,a_j,y_j,t_j,\sigma_j)$ 构成，其中 $o_j$ 为观测，$a_j$ 为动作或意图，$y_j$ 为结果，$t_j$ 为时间索引，$\sigma_j$ 为来源和置信元数据。结构记忆 $\Theta$ 是支持预测或策略读出的慢变量。回放调度器以权重
\[
w_j\propto \alpha_n n_j+\alpha_r r_j+\alpha_v v_j+\alpha_d d_j
\]
选择事件，$n_j$ 表示新颖性，$r_j$ 表示任务相关性，$v_j$ 表示价值或情绪显著度，$d_j$ 表示与保留集的覆盖缺口；各项必须归一到可比较尺度。高权重只表示优先检查，不表示应当永久保存。回放产生候选梯度 $g_j$，聚合更新为 $\Delta\Theta=-\alpha\sum_j\bar w_j g_j$，随后必须经过冲突检测、留出验证与版本提交。原始事件仍保留在 $E$ 或外部记忆Me中，以便发现结构归纳错误后追溯。

这一过程对应三相记忆的明确分工。$h(t)$ 承载当前回放片段、来源标签和冲突信号；$E$ 保持带时间与情境的经历，允许多个互不一致的版本并存；$\Theta$ 只吸收跨事件稳定且通过验证的结构。回放不是把全部情节逐字复制进参数，而是让模型在受控条件下重新计算预测、价值和因果候选。若新经验只在单一场景成立，就应保持为情景索引或外部规则，不能因为反复出现便升级为普遍结构。遗忘也不是简单删除：可以降低检索优先级、压缩细节、保留摘要和证据指针，使未来任务仍能恢复来源。

情绪加工需要另行定义。令事件的价值估计为 $V_k(e_j)$、不确定度为 $u_k(e_j)$、身体或资源状态为 $b_k$。离线再估计可写为
\[
V_{k+1}(e_j)=V_k(e_j)+\beta\bigl(\widehat V(e_j\mid \Theta_k,b_k)-V_k(e_j)\bigr),
\]
其中 $\beta\in[0,1]$ 是更新率。该式表示在新上下文中调整评价，不表示恐惧“能量”因重复激活而被守恒地消耗。反复暴露可能降低、维持或增强反应，取决于安全线索、可控性、预测误差和个体历史；噩梦也可能重复强化威胁模型。因而对创伤、焦虑或睡眠障碍的解释与干预必须服从临床证据，本模型不能作为诊断工具，更不能据此自行安排高强度重放。

工程上，AgentOS可在执行器隔离的沙盒中回放失败轨迹：例如机器人白天抓取易碎物失败，夜间维护从E提取传感器、动作和结果，使用模拟器生成摩擦系数与视角变体，再对候选策略做离线训练。若新策略只在模拟器内改善而真实保留集下降，就拒绝提交；若通过，则以新版本写入$\Theta$，并保留失败轨迹和测试证据。对于涉及身份、权限或高风险目标的记录，Boundary还应限制生成器访问范围，防止维护过程泄露隐私或复现危险动作。对其他架构，同一原则可由数据库事务、持续学习回放或模型注册表实现，不要求使用三相模块名称。

验收不仅看次日是否“记得”，还要分解为事实保持、技能迁移、来源准确、旧任务回归、情绪校准与副作用。对人类研究，梦的报告可作为体验证据之一，但不能单独定位回放内容；对机器系统，完整日志能证明哪些样本被使用，却仍不能证明更新获得了正确因果结构。我们保留原章“平滑与刻蚀同时发生”的问题意识，把它改写为压缩、保留、重估与结构提交之间的多目标权衡：只有可追溯的候选更新和独立验证，才能把离线活动与长期能力变化联系起来。

\section{噩梦、清明梦与睡眠瘫痪的混合状态}
\label{sec:section_4147}
噩梦、清明梦和睡眠瘫痪不宜排列成同一“能量密度”轴上的三个相态。它们涉及威胁内容、元认知、觉醒水平、运动抑制、外部感知和记忆报告等不同维度，并且可以部分共现。我们定义模式向量 $q_k=(A_k,M_k,G^a_k,G^e_k,V_k)$：$A_k$ 是觉醒水平，$M_k$ 是元认知增益，$G^a_k$ 是动作通道增益，$G^e_k$ 是外部感知增益，$V_k$ 是当前威胁评价。各量须由具体实验指标估计，不能把它们当作脑内直接可读的统一坐标。

噩梦可建模为威胁内容与唤醒反馈形成的高驻留区域。令威胁状态占用为 $r_k\in[0,1]$，其更新受记忆触发、身体信号、控制增益和随机扰动共同影响：
\[
r_{k+1}=\sigma(\alpha r_k+\beta V_k+\gamma b_k-\delta M_k+\xi_k).
\]
当自反馈 $\alpha$ 较大且纠偏项不足时，系统可能在多个步长内持续生成威胁片段；但有界激活函数意味着这不是数学发散，也不能据此称为无散旋涡。惊醒可以由威胁评价、内感受异常或外部刺激共同触发。阈值模型只能预测觉醒概率随变量变化，不能证明某个“惊奇激波”冲破睡眠。对反复噩梦，应记录频率、痛苦程度和日间功能，并交由合适的临床框架评估。

清明梦体现的是 $M_k$ 上升而睡眠相关的输入与动作配置尚未完全恢复。梦者可能识别状态来源、保持意图并影响内容，但控制范围有强弱差异，也可能迅速失去清明或醒来。我们的读出条件是：当来源判断对内部生成给出高置信标签，且自我状态报告在连续窗口内一致，才记为元认知命中；场景可控度则另行测量。这样可以避免把“知道在做梦”“能改变梦境”和“拥有外界行动能力”混为一谈。即便内容响应意图，也只说明生成闭环改变，不意味着物理法则暂停。

睡眠瘫痪则是觉醒与运动门控不同步的候选描述：$A_k$ 和部分感知、元认知已经提高，而 $G^a_k$ 仍低。主体发出动作意图却得不到预期本体反馈，形成显著预测误差
\[
\varepsilon_k^{act}=\widehat y_k^{body}-y_k^{body}.
\]
该误差可以与呼吸感受、姿态、残余梦像和威胁先验共同进入解释模型，产生“房间里有人”“胸口受压”或黑影等体验。这里不存在必须守恒并反弹的意图能量；运动命令未执行并不会自动变成反向激波。所谓“鬼”的形象可作为大脑对不完整感觉进行解释性生成的一种可能，但文化经验、期待和个体差异都会改变内容，不能声称单一机制已经解释全部案例。

这一分解也给出反例。高威胁梦不一定导致惊醒，清明梦不一定完全摆脱恐惧，睡眠瘫痪可以没有视觉幻觉，短暂无法行动也不必然属于睡眠瘫痪。模型若把这些差异都吸收到不可观测的“能量”或“几何张力”，就失去可证伪性。可行研究应同步记录睡眠阶段、生理信号、唤醒时报告、动作尝试与后续回忆，比较不同解释对数据的预测。报告中的时间顺序也要谨慎，因为醒后叙述可能把先后发生的感觉与解释重排。

AgentOS可以借这些现象设计故障测试，却不能将其冒充人类机制。系统应模拟“计划器恢复但执行器仍锁定”“威胁评估高但外部证据低”“来源标签恢复但生成器仍运行”等混合状态，检查是否出现重试风暴、虚假入侵报警或未经授权的动作。TCE负责模式迁移，Boundary保持硬权限，CCM限制重试负荷，SPC回报当前自我与来源状态；替代架构也可用状态机和安全监控器实现。通过这些测试，我们得到的是异步恢复的工程鲁棒性，而不是对噩梦或睡眠瘫痪的医学证明。

\section{机器离线维护：回放、反事实与安全提交}
\label{sec:robot_dreaming}
“机器人必须做梦吗”需要拆成两个问题：长期智能系统是否需要离线维护，以及这种维护是否应被称为梦。前者常有工程价值，后者只是功能类比。具备只读知识库、冻结模型或外部训练流水线的系统，可能无需自身生成梦样体验；连续学习的具身Agent则可能受益于回放、反事实模拟、模型压缩和安全回归。我们因此不宣称做梦是所有AGI的热力学必需条件，而给出一套可选的维护协议及其适用边界。

维护周期从一致快照开始。系统冻结可变结构的版本 $\Theta_k$，记录外部记忆Me、情景记忆E、工具接口和环境模型的版本，随后把真实执行器置为只读。候选数据分为失败回放、代表性成功样本、生成变体和极端反事实；每条生成数据必须携带谱系，注明生成器、种子、父样本和约束。训练器输出候选 $\Theta'_k$，而不是直接覆盖生产模型。评估器在隔离环境中比较能力、校准、安全、资源和旧任务回归，只有满足联合接收谓词
\[
\mathsf{Accept}(\Theta'_k)=
\bigwedge_{j=1}^{m}\left(Q_j(\Theta'_k)-Q_j(\Theta_k)\ge -\epsilon_j\right)
\land \mathsf{Safe}(\Theta'_k)
\]
时才允许提交。阈值 $\epsilon_j$ 由任务风险给定；接收谓词是工程规则，不是智能的一般定律。

所谓灾难性遗忘在这里表现为旧任务或旧分布性能显著下降，而非认知流形曲率变尖。缓解方法包括混合回放、参数约束、模块隔离、蒸馏、外部知识保留和版本回滚，各自有容量与计算代价。生成式扩散可以扩展样本覆盖，也会复制偏见或制造不可实现状态；低学习率可减少单步漂移，却不能消除长期累积误差。维护调度器应估计预计收益、验证成本、停机窗口和真实能耗，必要时把训练卸载到外部计算资源。这里的实际电耗可以用焦耳或千瓦时测量，训练损失和采样熵则是不同量。

大语言模型在事实任务中产生幻觉，也不应直接称为“清醒做梦”。它通常意味着生成分布、上下文证据、工具调用、训练目标和答案校准之间出现错配。创造性写作中的虚构并非故障，事实问答中的无根据断言才需要来源核验。AgentOS实例可让Boundary根据任务类型设定证据门槛：事实模式要求Me检索、引用与不确定性报告；创意模式允许更宽的生成跨度，但输出仍标注为构造内容。SPC压缩“我掌握何种证据”的自我状态，TCE调节探索与核验，CCM监控上下文和工具预算。其他系统完全可以用检索增强、验证器或拒答策略达到同一目的。

以火星建筑机器人为例，白天系统在低重力、粉尘和通信延迟下执行施工，E记录材料断裂、工具打滑和调度冲突。离线阶段可在数字环境中回放失败，改变土壤参数、队伍分工和设备故障组合，寻找风险较低的计划；但仿真胜出不等于现场成功。候选策略必须经过硬件在环测试、资源预算、通信中断和安全停机验证，之后才能逐级部署。若环境模型版本过旧，Boundary应阻止提交并要求刷新外部数据。这个案例保留原章“安全沙盒训练极端情形”的意图，同时明确模拟与现实之间的验证桥梁。

完整协议还包括退出条件。维护超时、生成样本比例过高、回归测试失败、签名不匹配或外部告警到达时，系统应中止训练，恢复最后可信版本，并把未完成候选隔离。连续多轮无收益说明维护目标或数据选择需要修订，而不是无限增加随机性。验收报告应列出新旧任务矩阵、最坏下降、校准误差、危险动作率、回滚时间、数据谱系完整性和实际能耗。由此，“机器之梦”成为一个受权限、证据和版本治理约束的离线生成协议；它可以服务于长期适应，但既不证明机器具有与人类相同的梦境体验，也不构成所有智能系统唯一的自维护方式。
